\section{Step 86 Retain mode, strides and success count}
Fill \texttt{pack\_state}. Retain scalar/MSA mode in \texttt{\$v1} before reusing \texttt{\$t9} for destination stride in bytes. Retain the active success count $4wh$ in \texttt{\$t8}. Then arrange the loop arguments: \texttt{\$a0} remains at the source selector, \texttt{\$a1} becomes the current destination row start, \texttt{\$a2} becomes width, and \texttt{\$a3} becomes rows remaining.

The worked direct call retains stride 68 and active count 40; the file caller retains stride 2048 and the same active count. Keep the full-width destination pointer when moving it between registers. A valid address can exceed 4 GiB even though the image dimensions are small.

The supplied branch bypasses vector-mask loading in scalar mode. In MSA mode, mask construction and loading are still unfinished. A temporary development return is expected until those later regions are completed; it is distinct from the format or capacity failure result.

\paragraph{Observe} Use \texttt{s86-1}; stop before \texttt{pack\_state\_ready}. Inspect the state described above before continuing.

\Needspace{7\baselineskip}
\paragraph{Check} Which count belongs in the success result?
\begin{itemize}
  \item[A.] The destination allocation capacity.
  \item[B.] The compressed file length.
  \item[C.] The active pixel-byte count.
\end{itemize}
